\documentclass[11pt]{article}
\usepackage{mathblatt}
\makeatletter
\setlength{\mbpfbe}{0mm}% keine Punkte (Bauregel 6.4)
\makeatother
% Rechenraum als Karo (Bauregel 6.9): n Rechenschritte = n mal 10 mm
\newcommand{\karo}[1]{\par\smallskip\noindent\begin{tikzpicture}
  \clip (0,0) rectangle ({\linewidth-1mm},{#1*10mm});
  \draw[black!16,line width=0.3pt,step=5mm] (0,0) grid ({\linewidth},{#1*10mm});
  \end{tikzpicture}\par}
\newcommand{\linie}[1]{\,{\color{black!45}\rule[-1pt]{#1}{0.4pt}}\,}
\newcommand{\pkt}{\textperiodcentered{}}
\begin{document}

\pfheftstil
\pflaufkopf[P10 \pkt{} mehr]{Winkel ohne Rechnung bestimmen oder begründen \pkt{} Lösungen}{PG10-4}\pfseitevon
\pfheftkopf{Winkel ohne Rechnung bestimmen oder begründen}{Lösungen \pkt{} mehr}
\begin{pfloesungb}{}
\lzb{1.}{}{$\alpha$ = 35° (parallele Seiten oben und unten)}{$\alpha$ = 180° $-$ 145° = 35° (parallele Seiten oben und unten)}
\end{pfloesungb}
\begin{pfloesungb}{}
\lzb{2.}{}{$\gamma$ = 79,5°}{$\gamma$ = 180° $-$ 64,4° $-$ 36,1° = 79,5°}
\end{pfloesungb}
\begin{pfloesungb}{}
\lzb{3.}{}{$\delta$ = 28°}{$\delta$ = 180° $-$ 118° $-$ 34° = 28°}
\end{pfloesungb}
\begin{pfloesungb}{}
\lzb{4.}{}{$\delta$ = 56°}{$\sphericalangle$ACB = 90° $-$ 30° = 60° $\Rightarrow$ $\sphericalangle$ACD = 180° $-$ 60° = 120° $\Rightarrow$ $\delta$ = 180° $-$ 120° $-$ 4° = 56°}
\end{pfloesungb}
\end{document}